\documentclass[11pt]{article}

\usepackage[a4paper,margin=28mm,headheight=15pt]{geometry}
\usepackage[T1]{fontenc}
\usepackage[utf8]{inputenc}
\usepackage{lmodern}
\usepackage{microtype}
\usepackage{amsmath,amssymb,amsthm,mathtools,mathrsfs}
\usepackage{booktabs,longtable,array,tabularx}
\usepackage{enumitem}
\usepackage{xcolor}
\usepackage{hyperref}
\usepackage{url}
\usepackage{listings}
\usepackage{fancyhdr}
\usepackage{lastpage}
\usepackage{graphicx}
\usepackage{caption}

\definecolor{statusblue}{RGB}{28,74,125}
\definecolor{statusgreen}{RGB}{30,105,65}
\definecolor{statusorange}{RGB}{150,82,0}
\definecolor{statusred}{RGB}{145,35,35}
\definecolor{lightgray}{RGB}{246,246,246}

\hypersetup{
  colorlinks=true,
  linkcolor=statusblue,
  citecolor=statusgreen,
  urlcolor=statusblue,
  pdftitle={Finite-Field Rank Tomography and Local--Global Rigidity for the Alaoglu--Erdos Conjecture},
  pdfauthor={OpenAI GPT-5.6 Pro},
  pdfsubject={Nonduplicative progress report on the Alaoglu--Erdos conjecture}
}

\pagestyle{fancy}
\fancyhf{}
\lhead{Alaoglu--Erd\H{o}s: finite-field rank tomography}
\rhead{22 August 2026}
\cfoot{\thepage\ of \pageref{LastPage}}

\newtheorem{theorem}{Theorem}[section]
\newtheorem{proposition}[theorem]{Proposition}
\newtheorem{lemma}[theorem]{Lemma}
\newtheorem{corollary}[theorem]{Corollary}
\newtheorem{conjecture}[theorem]{Conjecture}
\newtheorem{question}[theorem]{Question}
\theoremstyle{definition}
\newtheorem{definition}[theorem]{Definition}
\newtheorem{algorithm}[theorem]{Algorithm}
\theoremstyle{remark}
\newtheorem{remark}[theorem]{Remark}

\newcommand{\Z}{\mathbb Z}
\newcommand{\Q}{\mathbb Q}
\newcommand{\R}{\mathbb R}
\newcommand{\F}{\mathbb F}
\newcommand{\Gm}{\mathbb G_m}
\newcommand{\ord}{\operatorname{ord}}
\newcommand{\im}{\operatorname{im}}
\newcommand{\rank}{\operatorname{rank}}
\newcommand{\rad}{\operatorname{rad}}
\newcommand{\vp}{v_p}
\newcommand{\crow}{c^{\mathrm{row}}}
\newcommand{\ccol}{c^{\mathrm{col}}}
\newcommand{\E}{\mathcal E}
\newcommand{\Torus}{T}
\newcommand{\Nzero}{\mathbb N_0}
\newcommand{\lcm}{\operatorname{lcm}}
\newcommand{\diag}{\operatorname{diag}}
\newenvironment{resultbox}{\begin{quote}\small}{\end{quote}}
\newcommand{\statusproved}{\textcolor{statusgreen}{\textbf{PROVED HERE}}}
\newcommand{\statusknown}{\textcolor{statusblue}{\textbf{KNOWN / IMPORTED}}}
\newcommand{\statuscompute}{\textcolor{statusorange}{\textbf{COMPUTATIONAL}}}
\newcommand{\statusopen}{\textcolor{statusred}{\textbf{OPEN GAP}}}

\setlist[itemize]{topsep=4pt,itemsep=2pt,parsep=0pt}
\setlist[enumerate]{topsep=4pt,itemsep=2pt,parsep=0pt}
\setlength{\parindent}{0pt}
\setlength{\parskip}{5pt}
\allowdisplaybreaks

\lstset{
  basicstyle=\ttfamily\scriptsize,
  columns=fullflexible,
  breaklines=true,
  frame=single,
  backgroundcolor=\color{lightgray},
  showstringspaces=false,
  keepspaces=true,
  tabsize=4,
  numbers=left,
  numberstyle=\tiny,
  xleftmargin=2em,
  framexleftmargin=1.5em
}

\title{\bfseries Finite-Field Rank Tomography and Local--Global Rigidity\\[4pt]
for the Alaoglu--Erd\H{o}s Conjecture\\[10pt]
\large A nonduplicative progress report}
\author{OpenAI GPT-5.6 Pro}
\date{22 August 2026}

\begin{document}
\maketitle

\begin{abstract}
Let $x\in\R$ and suppose that $2^x$ and $3^x$ are integers.  The Alaoglu--Erd\H{o}s conjecture asserts that $x$ is an integer.  The attached unified report already develops a very large collection of analytic, algebraic, determinant, $p$-adic, interpolation, recurrence, and dynamical approaches.  The present report deliberately does not recapitulate that material.  Instead it introduces a pair-specific finite-place formalism for a hypothetical counterexample
\[
P=(2,3),\qquad Q=(M,A)=(2^\beta,3^\beta)\in(\Q^\times)^2,
\]
where $\beta\notin\Z$ is positive and $M,A\in\Z$.

The main new result is an exact \emph{Smith tomography inequality}.  For every good prime power $p^t$, the number of distinct reductions of a chosen row support under $(a,b)\mapsto P^aQ^b$, and the dual number of character signatures of a chosen column support, bound every layer of the Smith normal form of the associated integer character-evaluation determinant.  Short vectors in the resulting congruence lattices yield weighted common-divisor certificates.  A discrete-logarithm matrix $B_p$ makes the row and column pictures transposes of one another and gives an exact formula for their full finite-field image sizes.  These ingredients lead to a finite, exact, and readily formalizable contradiction certificate: a computable orbit-collapse energy must exceed an explicit archimedean Hadamard budget.

Two local--global consequences are also isolated.  Perucca's cyclic-subgroup detection theorem implies that almost-everywhere local membership $\bar Q_p\in\langle\bar P_p\rangle$ would already prove the conjecture.  A support theorem for reduction orders likewise forces integrality if $\ord(\bar Q_p)$ divides $\ord(\bar P_p)$ for almost all primes.  Conversely, fixed local slopes are rigid, so any useful slopes in a counterexample must drift.  A quantitative no-go theorem shows that cyclic image cardinality alone is far too weak; even a weighted certificate built from all short local cyclic slopes cannot close the present Hadamard estimate throughout a substantial range of $\beta$.  Thus the remaining bridge is sharply identified: one needs many moving short slopes, exceptional $p$-adic depth, a stronger archimedean determinant bound, or a hybrid of these.

Exact computations for primes below $5000$ compare a genuine integral-power pair, two generic pairs, and an archimedean near-miss with logarithmic discrepancy about $1.36\times10^{-9}$.  The integral pair collapses at every tested prime, whereas the near-miss remains finite-field generic at almost every prime.  No complete proof of the conjecture is claimed.  The report's contribution is a new exact framework, several rigorous reduction and obstruction theorems, reproducible experiments, and a narrowly specified next proof target.
\end{abstract}

\tableofcontents
\newpage

\section{Scope, status, and nonduplication}

\subsection{What is and is not being claimed}

The status of the central problem is:
\[
\boxed{\text{The Alaoglu--Erd\H{o}s conjecture is not proved in this report.}}
\]
That sentence is not a disclaimer replacing mathematical work.  It is a precise boundary around the work below.  The report proves new finite-field determinant inequalities, derives exact local--global reductions, identifies quantitative ceilings for tempting strategies, and supplies a reproducible computation.  The final implication from the real proportionality
\[
(\log M,\log A)=\beta(\log2,\log3)
\]
to sufficiently abundant finite-place orbit collapse remains open.

Every principal statement is marked conceptually as one of the following.
\begin{itemize}
\item \statusproved: proved in this report.
\item \statusknown: a standard theorem or a fact already present in the attached report, imported only as needed.
\item \statuscompute: exact finite computation, not a theorem about all primes.
\item \statusopen: a sharply formulated missing bridge or conditional target.
\end{itemize}

\subsection{Nonduplication ledger}

The attached manuscript is more than twenty thousand lines long and contains over one hundred sections.  It already treats the minimal-counterexample normalization, total positivity of exponential kernels, generalized Vandermonde determinants, universal modular Smith divisors, finite-place and Iwasawa-log methods, gcd bounds, recurrence sequences, interpolation determinants, Kummer theory, Loewner matrices, canonical products, dynamics, entropy, and the restricted Matrix Coefficient Conjecture.  Repeating those chapters would obscure rather than advance the problem.

\begin{longtable}{>{\raggedright\arraybackslash}p{0.26\textwidth} >{\raggedright\arraybackslash}p{0.31\textwidth} >{\raggedright\arraybackslash}p{0.35\textwidth}}
\toprule
Topic & Treatment here & Relation to the attached report\\
\midrule
Counterexample notation & One compact section & Imported baseline only\\
Exponential determinant & One proposition and one Hadamard bound & Imported baseline only\\
Universal $d$-variable Smith divisor & Quoted for scale comparison & Explicitly not claimed as new\\
Exact row orbit map modulo $p^t$ & Full development & New pair-specific object\\
Exact dual column-signature map & Full development & New pair-specific object\\
Smith tomography by orbit counts & Theorem with proof & New\\
Short relation amplification & Theorem with weighted gcd corollaries & New\\
Discrete-log image formula and transpose duality & Theorem with proof & New\\
Full character-table rank criterion & Proposition with proof & New\\
Perucca support theorems specialized to $(2,3),(M,A)$ & Detailed consequences & New specialization\\
Cyclic-cardinality ceiling and short-slope ceiling & Asymptotic no-go theorems & New\\
Prime computation below $5000$ & Exact reproducible scout & New\\
\bottomrule
\end{longtable}

\subsection{Public-record audit of the claimed rival breakthroughs}

Because the prompt explicitly invokes very recent claims, a source audit was performed.  This audit is not a mathematical peer review and does not attempt to adjudicate unpublished work.

\begin{enumerate}
\item Anthropic publicly released Claude Fable 5 on 11 August 2026 and published a system card describing several research collaborations \cite{AnthropicFable,AnthropicSystemCard}.  Thus the model name is not fictitious.
\item A public preprint by Alp\"oge gives explicit counterexamples to the Jacobian conjecture in every dimension $n\ge3$; the public project page describes Fable's role in the computational search and proof development \cite{AlpogeJacobian,ByClaudeJacobian}.  This refutes the conjecture in those dimensions; it does not settle the still-distinct plane case $n=2$.
\item Public project pages report Hadamard constructions including order $668$, and one page claims completion of a larger list of previously open orders below $2000$ \cite{ByClaudeHadamard,VibeHadamard}.  The public benchmark pages were not fully synchronized when checked, so this report does not certify the exact count stated in the prompt.
\item No citable public construction of an elliptic curve over $\Q$ of algebraic rank at least $30$ was located.  The public benchmark still described rank $29$ as the record and rank $30$ as the open target \cite{VibeRank30}.  A private or not-yet-indexed result cannot be ruled out, but it cannot be used as evidence here.
\item No public proof or Lean development of the Alaoglu--Erd\H{o}s conjecture was located.  The rival team's private message is therefore not independently verifiable.
\end{enumerate}

The appropriate response to competitive pressure is not to lower the standard of proof.  It is to extract a mathematically exact target that can be attacked, checked, and formalized.

\section{Minimal inherited framework}\label{sec:grid}

\subsection{Normalization of a hypothetical counterexample}

\begin{conjecture}[Alaoglu--Erd\H{o}s]
If $x\in\R$ and $2^x,3^x\in\Z$, then $x\in\Z$.
\end{conjecture}

Negative $x$ is impossible unless $2^x$ and $3^x$ are nonzero integers of absolute value below $1$.  Hence any nontrivial counterexample is positive.  Assume for the rest of the report that
\[
\beta>0,\qquad \beta\notin\Z,\qquad M=2^\beta\in\Z,\qquad A=3^\beta\in\Z.
\]
The earlier report proves much stronger finite exclusions; for orientation, one may keep in mind the established lower range $\beta>27$.

Set
\[
\Torus=(\Gm)^2,\qquad P=(2,3),\qquad Q=(M,A)\in \Torus(\Q).
\]
The real logarithm matrix
\[
\Lambda=\begin{pmatrix}
\log2&\log M\\
\log3&\log A
\end{pmatrix}
=\begin{pmatrix}
\log2&\beta\log2\\
\log3&\beta\log3
\end{pmatrix}
\]
has rank one.  The finite-place program developed below asks whether this exact real rank-one relation must cast a quantitatively strong ``shadow'' modulo many prime powers.

\begin{lemma}[Irrationality and transcendence of $\beta$]\label{lem:beta-trans}
A nonintegral $\beta$ satisfying the displayed hypotheses is irrational and transcendental.
\end{lemma}

\begin{proof}
If $\beta=a/b\in\Q$ in lowest terms with $b>0$, then $M^b=2^a$.  Unique factorization forces $b\mid a$, hence $b=1$ and $\beta\in\Z$, a contradiction.  Thus $\beta$ is irrational.  If it were algebraic, then it would be algebraic irrational and the Gelfond--Schneider theorem would make $2^\beta$ transcendental, contradicting $M\in\Z$.
\end{proof}

\begin{lemma}[Global independence]\label{lem:global-independence}
The points $P$ and $Q$ generate a free rank-two subgroup of $\Torus(\Q)$: if $P^uQ^v=(1,1)$ for $u,v\in\Z$, then $u=v=0$.
\end{lemma}

\begin{proof}
The first coordinate gives $2^uM^v=2^{u+v\beta}=1$, so $u+v\beta=0$.  Lemma~\ref{lem:beta-trans} makes $\beta$ irrational, hence $u=v=0$.
\end{proof}

\subsection{The integer character grid}

For finite supports
\[
\Omega\subset\Z_{\ge0}^2,\qquad \Sigma\subset\Z_{\ge0}^2,
\]
define the character-evaluation matrix
\begin{equation}\label{eq:evaluation-matrix}
E_{\Omega,\Sigma}[(a,b),(r,s)]
=\chi_{r,s}(P^aQ^b)
=(2^aM^b)^r(3^aA^b)^s.
\end{equation}
Every entry is an integer.  If $|\Omega|=|\Sigma|=m$, write
\[
D_{\Omega,\Sigma}=\det E_{\Omega,\Sigma}\in\Z.
\]
Using $M=2^\beta$ and $A=3^\beta$, the same entry is
\begin{equation}\label{eq:exp-kernel}
E_{\Omega,\Sigma}[(a,b),(r,s)]
=\exp\!\left((a+b\beta)(r\log2+s\log3)\right).
\end{equation}

\begin{proposition}[Nonvanishing baseline; \statusknown]\label{prop:nonvanishing}
If the values $a+b\beta$ are distinct on $\Omega$ and the values $r\log2+s\log3$ are distinct on $\Sigma$, then $D_{\Omega,\Sigma}\ne0$.
\end{proposition}

\begin{proof}
After sorting both parameter sets increasingly, the matrix in \eqref{eq:exp-kernel} is a minor of the strictly totally positive kernel $K(x,y)=e^{xy}$.  Equivalently, a nonzero exponential polynomial $\sum_{j=1}^m c_j e^{\lambda_j x}$ with distinct real $\lambda_j$ has at most $m-1$ real zeros counted with multiplicity; induction after applying $e^{-\lambda_1x}\frac{d}{dx}$ proves the Chebyshev property.  Hence the interpolation matrix is nonsingular.
\end{proof}

For rectangular boxes
\begin{equation}\label{eq:boxes}
\Omega_{U,V}=[0,U)\times[0,V),\qquad
\Sigma_{R,S}=[0,R)\times[0,S),
\end{equation}
with $UV=RS=m$, distinctness follows from the irrationality of $\beta$ and of $\log_2 3$.  Thus $|D_{\Omega,
\Sigma}|\ge1$.  The strategy is to prove that the same nonzero integer is divisible by more prime powers than its archimedean size allows.

\section{Exact finite-place orbit maps}

\subsection{Good primes and the two reduction maps}

A prime $p$ is called \emph{good} if
\[
p\nmid6MA.
\]
For a good prime and $t\ge1$, all four coordinates are units modulo $p^t$.  Define the row orbit map
\begin{equation}\label{eq:row-map}
\iota_{p^t}:\Z^2\longrightarrow \bigl((\Z/p^t\Z)^\times\bigr)^2,
\qquad
(a,b)\longmapsto (2^aM^b,3^aA^b)\pmod{p^t}.
\end{equation}
For a finite row support $\Omega$, put
\begin{equation}\label{eq:crow}
\crow_{p,t}(\Omega)=|\iota_{p^t}(\Omega)|.
\end{equation}

The dual column-signature map is
\begin{equation}\label{eq:col-map}
\iota^*_{p^t}:\Z^2\longrightarrow \bigl((\Z/p^t\Z)^\times\bigr)^2,
\qquad
(r,s)\longmapsto (2^r3^s,M^rA^s)\pmod{p^t},
\end{equation}
with
\begin{equation}\label{eq:ccol}
\ccol_{p,t}(\Sigma)=|\iota^*_{p^t}(\Sigma)|.
\end{equation}
The row map records points of the subgroup generated by $P$ and $Q$.  The column map records how ambient monomial characters restrict to that subgroup.

\subsection{Exact congruence lattices}

The kernels are rank-two lattices of finite index:
\begin{align}
L^{\mathrm{row}}_{p^t}
&=\{(h,k)\in\Z^2:2^hM^k\equiv1,\ 3^hA^k\equiv1\pmod{p^t}\},\label{eq:row-lattice}\\
L^{\mathrm{col}}_{p^t}
&=\{(r,s)\in\Z^2:2^r3^s\equiv1,\ M^rA^s\equiv1\pmod{p^t}\}.\label{eq:col-lattice}
\end{align}
Consequently,
\begin{equation}\label{eq:box-class-count}
\crow_{p,t}(\Omega)=|\Omega/L^{\mathrm{row}}_{p^t}|,
\qquad
\ccol_{p,t}(\Sigma)=|\Sigma/L^{\mathrm{col}}_{p^t}|,
\end{equation}
where the right-hand side means the number of lattice cosets met by the finite set.  This is an exact pair-specific object, not a universal bound obtained merely by counting all residues.

\begin{remark}
The old universal argument only uses
\[
\crow_{p,t}(\Omega),\ccol_{p,t}(\Sigma)\le\varphi(p^t)^2<p^{2t}.
\]
The new program computes or bounds the actual box sections of the two lattices.  A short lattice vector can make the count dramatically smaller even when the full finite subgroup is large.
\end{remark}

\section{Pair-specific Smith tomography}

\subsection{Smith layers as generator counts}

Let $E=E_{\Omega,\Sigma}$ be an $m\times m$ integer matrix with nonzero determinant $D$.  Fix a prime $p$.  Its Smith normal form over $\Z_p$ has diagonal entries
\[
p^{\nu_1},\ldots,p^{\nu_m},
\qquad 0\le\nu_1\le\cdots\le\nu_m,
\]
up to multiplication by $p$-adic units.  For $t\ge1$ define
\begin{equation}\label{eq:smith-rank-layer}
r_{p,t}=\#\{i:\nu_i<t\}.
\end{equation}
The elementary Ferrers-diagram identity is
\begin{equation}\label{eq:ferrers}
\vp(D)=\sum_{i=1}^m\nu_i
=\sum_{t\ge1}(m-r_{p,t}).
\end{equation}
Only finitely many summands are nonzero.

\begin{lemma}[Module interpretation]\label{lem:module-generators}
The integer $r_{p,t}$ is the minimal number of generators of the row module of $E$ over $\Z/p^t\Z$.  It is also the minimal number of generators of the column module.
\end{lemma}

\begin{proof}
Invertible row and column operations do not change either minimal generator number.  For the Smith diagonal, the row module is
\[
\bigoplus_{\nu_i<t}p^{\nu_i}(\Z/p^t\Z).
\]
The ring $\Z/p^t\Z$ is local.  Modulo its maximal ideal, each nonzero summand contributes one copy of $\F_p$, so Nakayama's lemma gives exactly one required generator per index with $\nu_i<t$.  The column argument is identical.
\end{proof}

\subsection{The tomography inequality}

\begin{theorem}[Finite-field Smith tomography; \statusproved]\label{thm:smith-tomography}
Let $p$ be good, let $t\ge1$, and let $E_{\Omega,\Sigma}$ be square of size $m$ with nonzero determinant.  Then
\begin{equation}\label{eq:rank-tomography}
r_{p,t}\le
\min\bigl(\crow_{p,t}(\Omega),\ccol_{p,t}(\Sigma)\bigr).
\end{equation}
Consequently,
\begin{equation}\label{eq:valuation-tomography}
\vp(D_{\Omega,\Sigma})
\ge
\sum_{t\ge1}
\left(m-
\min\bigl(\crow_{p,t}(\Omega),\ccol_{p,t}(\Sigma)\bigr)
\right).
\end{equation}
\end{theorem}

\begin{proof}
Reduce the matrix modulo $p^t$.  If two row indices $(a,b),(a',b')$ have the same image under $\iota_{p^t}$, then
\[
(2^aM^b,3^aA^b)\equiv
(2^{a'}M^{b'},3^{a'}A^{b'})\pmod{p^t}.
\]
Every monomial character therefore takes the same value on those two points, so the corresponding matrix rows are identical modulo $p^t$.  The row module is generated by at most one representative from each orbit class met by $\Omega$; hence its minimal number of generators is at most $\crow_{p,t}(\Omega)$.  Lemma~\ref{lem:module-generators} gives
$r_{p,t}\le\crow_{p,t}(\Omega)$.

Likewise, if two columns have the same image under $\iota^*_{p^t}$, then their quotient character is trivial on both generators $P$ and $Q$, hence on every row point $P^aQ^b$.  The two columns are identical modulo $p^t$, and
$r_{p,t}\le\ccol_{p,t}(\Sigma)$.  This proves \eqref{eq:rank-tomography}.  Summing $m-r_{p,t}$ and using \eqref{eq:ferrers} proves \eqref{eq:valuation-tomography}.
\end{proof}

\begin{remark}[Why this is stronger than rank modulo $p$]
The ordinary rank defect modulo $p$ controls only $m-r_{p,1}$.  Theorem~\ref{thm:smith-tomography} controls every threshold $t$ and therefore every layer of the Smith profile.  A congruence persisting from $p$ to $p^t$ is counted $t$ times, exactly as its valuation demands.
\end{remark}

\begin{corollary}[Global orbit-collapse divisor]\label{cor:global-orbit-divisor}
For any finite set $\mathcal P$ of good primes,
\begin{equation}\label{eq:global-orbit-divisor}
\prod_{p\in\mathcal P}
 p^{\sum_{t\ge1}
\left(m-\min(\crow_{p,t}(\Omega),\ccol_{p,t}(\Sigma))\right)}
\;\Bigm|\; D_{\Omega,\Sigma}.
\end{equation}
\end{corollary}

\begin{proof}
Apply Theorem~\ref{thm:smith-tomography} prime by prime.  Distinct prime powers are coprime.
\end{proof}

\subsection{Recovery of the universal bound and the genuinely new term}

The old universal estimate is recovered by the trivial inequalities
\[
\crow_{p,t},\ccol_{p,t}\le p^{2t}.
\]
Thus one gets
\[
\vp(D)\ge\sum_{t\ge1}(m-p^{2t})_+,
\]
and, after summing over primes,
\[
\log|D|\ge \left(\frac23+o(1)\right)m^{3/2}.
\]
This $2/3$ law is already in the attached report and is quoted only for comparison.  The new information is the nonnegative excess
\begin{equation}\label{eq:pair-specific-excess}
\Delta_{p,t}
=p^{2t}-\min(\crow_{p,t},\ccol_{p,t}),
\end{equation}
whenever the actual pair occupies fewer residues than the ambient two-dimensional torus permits.  The central quantitative problem is whether the sum of these pair-specific excesses can be forced to reach the archimedean scale $m^2$.

\section{Short congruence relations and weighted divisors}\label{sec:gcd}

\subsection{A box-translation graph}

Let $\Omega_{U,V}$ be the row box in \eqref{eq:boxes}.  For a nonzero vector $u=(h,k)\in\Z^2$ with $|h|<U$ and $|k|<V$, join $z$ to $z+u$ whenever both vertices lie in the box.  Orient each edge in the $u$ direction after replacing $u$ by $-u$ if necessary.

\begin{lemma}[Translation forest]\label{lem:translation-forest}
The graph is a disjoint union of paths and has exactly
\begin{equation}\label{eq:edge-count}
e_u=(U-|h|)(V-|k|)
\end{equation}
edges.  Hence it has $UV-e_u$ connected components.
\end{lemma}

\begin{proof}
The overlap of the box with its translate by $-u$ has the stated cardinality, giving the edge count.  A directed cycle would imply $nu=0$ in $\Z^2$ for some $n>0$, impossible because $u\ne0$.  Each vertex has at most one outgoing and one incoming edge, so every component is a path.  A forest with $UV$ vertices and $e_u$ edges has $UV-e_u$ components.
\end{proof}

\subsection{Short relation amplification}

\begin{theorem}[Short row relation amplification; \statusproved]\label{thm:short-relation}
Let $p$ be good and suppose that
\begin{equation}\label{eq:short-row-relation}
P^hQ^k\equiv(1,1)\pmod{p^t}
\end{equation}
for a nonzero $(h,k)\in\Z^2$ with $|h|<U$ and $|k|<V$.  For any column support $\Sigma$ of size $UV$ for which the evaluation matrix is square and nonsingular,
\begin{equation}\label{eq:short-relation-valuation}
\vp(D_{\Omega_{U,V},\Sigma})
\ge t(U-|h|)(V-|k|).
\end{equation}
The analogous assertion holds with rows and columns interchanged and a short vector in $L^{\mathrm{col}}_{p^t}$.
\end{theorem}

\begin{proof}
Along every edge $z\to z+(h,k)$ of the translation forest, relation \eqref{eq:short-row-relation} makes the two row points congruent modulo $p^t$, hence their matrix rows are congruent.  The row orbit map is constant on every connected component.  Lemma~\ref{lem:translation-forest} gives
\[
\crow_{p,j}(\Omega_{U,V})\le UV-(U-|h|)(V-|k|)
\]
for every $1\le j\le t$.  Apply the row half of Theorem~\ref{thm:smith-tomography} at each level $j$ and sum.
\end{proof}

The result is stronger than saying that the determinant vanishes modulo $p^t$.  A single relation is replicated across every overlapping translate in the box, and each replication contributes an independent Smith defect.

\subsection{Weighted gcd certificates}

For $0\le h<U$ and $0\le k<V$, not both zero, define
\begin{equation}\label{eq:row-gcd}
G^{\mathrm{row}}_{h,k}
=\gcd(2^hM^k-1,\,3^hA^k-1).
\end{equation}
For $0\le r<R$, $0\le s<S$, not both zero, define
\begin{equation}\label{eq:col-gcd}
G^{\mathrm{col}}_{r,s}
=\gcd(2^r3^s-1,\,M^rA^s-1).
\end{equation}

\begin{corollary}[Weighted common divisors; \statusproved]\label{cor:weighted-gcd}
For a square nonsingular box matrix with $UV=RS=m$,
\begin{align}
\left(G^{\mathrm{row}}_{h,k}\right)^{(U-h)(V-k)}&\mid D,\label{eq:row-weighted-divisor}\\
\left(G^{\mathrm{col}}_{r,s}\right)^{(R-r)(S-s)}&\mid D.\label{eq:col-weighted-divisor}
\end{align}
Consequently, the least common multiple of all prime-power-weighted quantities on either line divides $D$.  A product is not generally valid because the same prime valuation may be counted by several relations.
\end{corollary}

\begin{proof}
If $p^t$ divides $G^{\mathrm{row}}_{h,k}$, then $P^hQ^k\equiv1\pmod{p^t}$.  The direct row-difference argument in Theorem~\ref{thm:short-relation} works even at the finitely many primes excluded from the group notation, because the two displayed congruences themselves are sufficient.  Thus
\[
\vp(D)\ge t(U-h)(V-k).
\]
Take $t=\vp(G^{\mathrm{row}}_{h,k})$ for every prime.  The column case is identical.
\end{proof}

\begin{remark}[Why the lcm matters]
The attached report correctly emphasizes that multiplying many gcd divisors can overcount a prime.  The exact safe aggregate is
\[
\lcm_{h,k}
\left(G^{\mathrm{row}}_{h,k}\right)^{(U-h)(V-k)},
\]
whose $p$-valuation is the maximum, not the sum, of the competing weighted valuations.  The full Smith tomography energy can be larger because several independent lattice directions may create more components than any single relation detects.
\end{remark}

\section{Discrete logarithms and transpose duality}

\subsection{The local orbit matrix}

Let $p$ be an odd good prime, put $n=p-1$, choose a primitive root $g$ modulo $p$, and write $\ell_p(u)\in\Z/n\Z$ for the discrete logarithm of $u$.  Define
\begin{equation}\label{eq:Bp}
B_p=
\begin{pmatrix}
\ell_p(2)&\ell_p(M)\\
\ell_p(3)&\ell_p(A)
\end{pmatrix}\pmod n.
\end{equation}
Under the identification $(\F_p^\times)^2\cong(\Z/n\Z)^2$, the row map is multiplication by $B_p$:
\[
(a,b)^T\longmapsto B_p(a,b)^T.
\]
The column-signature map is multiplication by the transpose:
\[
(r,s)^T\longmapsto B_p^T(r,s)^T.
\]
This simple transpose is the finite-field origin of the row--column symmetry in Theorem~\ref{thm:smith-tomography}.

\begin{theorem}[Exact image size; \statusproved]\label{thm:image-size}
Choose arbitrary integer lifts $b_{ij}$ of the entries of $B_p$ and set
\begin{equation}\label{eq:qp}
q_p=\gcd\bigl(n^2,\,nb_{11},\,nb_{12},\,nb_{21},\,nb_{22},\,\det B_p\bigr).
\end{equation}
Then
\begin{equation}\label{eq:image-formula}
|\im B_p|=|\im B_p^T|=\frac{n^2}{q_p}.
\end{equation}
The formula is independent of the chosen primitive root and of the integer lifts.
\end{theorem}

\begin{proof}
Let $L\subset\Z^2$ be the lattice generated by the columns of the $2\times4$ matrix
\[
\bigl[nI_2\mid B_p\bigr].
\]
The quotient $\Z^2/L$ is the cokernel of $B_p:(\Z/n\Z)^2\to(\Z/n\Z)^2$.  Its order is the index $[\Z^2:L]$.  For a full-rank integer matrix with two rows, that index is the gcd of its $2\times2$ minors.  Those minors are precisely the six integers in \eqref{eq:qp}, up to signs.  Therefore the cokernel has order $q_p$, and the image has order $n^2/q_p$.  Transposition leaves the list of entries and the determinant unchanged up to sign, so it has the same image size.  Since the result computes an intrinsic group cardinality, it is independent of all choices.
\end{proof}

\begin{corollary}[Rank-one-scale reduction]\label{cor:rank-one-scale}
If
\begin{equation}\label{eq:detzero}
\det B_p\equiv0\pmod{p-1},
\end{equation}
then
\[
|\im B_p|=|\im B_p^T|\le p-1.
\]
\end{corollary}

\begin{proof}
Every number in the gcd \eqref{eq:qp} is then divisible by $n=p-1$, so $q_p\ge n$.
\end{proof}

The terminology ``rank-one-scale'' is deliberate.  Over the composite ring $\Z/(p-1)\Z$, determinant zero need not mean that the image is a free cyclic module, but its cardinality is at most the one-dimensional ambient scale $p-1$.  If $Q=P^c$ globally, then the second column of $B_p$ is $c$ times the first for every good prime, so \eqref{eq:detzero} holds everywhere.

\subsection{Exact rank when full local groups are sampled}

Let $H_p=\langle\bar P_p,\bar Q_p\rangle\subset(\F_p^\times)^2$.  The ambient monomial characters restrict to characters of $H_p$.

\begin{proposition}[Full character-table rank; \statusproved]\label{prop:character-rank}
Suppose the row support maps onto all of $H_p$ and the column support maps onto all distinct restricted characters of $H_p$.  Then the reduction of $E_{\Omega,\Sigma}$ modulo $p$ has rank exactly
\[
|H_p|=|\im B_p|.
\]
\end{proposition}

\begin{proof}
The image of $B_p^T$ has the same size as $H_p$ by Theorem~\ref{thm:image-size}; it is exactly the set of restriction signatures of ambient characters.  Choose one row representative for each $h\in H_p$ and one column representative for each restricted character $\chi\in\widehat H_p$.  The resulting square matrix is the character table $(\chi(h))$.  Its rows are orthogonal: for a nontrivial character $\psi$, choosing $h_0$ with $\psi(h_0)\ne1$ gives
\[
\sum_{h\in H_p}\psi(h)
=\sum_{h\in H_p}\psi(h_0h)
=\psi(h_0)\sum_{h\in H_p}\psi(h),
\]
so the sum vanishes.  Since $p\nmid|H_p|$, the character table is invertible over $\F_p$.  Duplicating rows or columns cannot increase rank, proving the claim.
\end{proof}

\begin{remark}
Proposition~\ref{prop:character-rank} distinguishes two causes of small determinant rank: a genuinely small subgroup $H_p$, and failure of the finite boxes to cover the available group or dual group.  For the determinant method, the second phenomenon is often more important.  A huge cyclic subgroup can still produce many box collisions when its local slope is short.
\end{remark}

\section{Local slopes: rigidity, divisibility, and the need to drift}
\label{sec:slopes}

The congruence lattice becomes especially concrete when the reduction of $Q$ lies in the cyclic subgroup generated by the reduction of $P$.  This is the finite-place analogue of a rational slope.  It is important to distinguish three assertions:
\begin{enumerate}[label=\textup{(\alph*)}]
\item the subgroup $\langle\bar P_p,\bar Q_p\rangle$ has small cardinality;
\item $\bar Q_p$ lies in $\langle\bar P_p\rangle$;
\item $\bar Q_p=\bar P_p^c$ for a \emph{small} integer representative $c$.
\end{enumerate}
The determinant sees increasingly more information as one passes from (a) to (c).  A cyclic image can still have order about $p$; local membership can require a slope of size about $p$; only a short slope gives a large collision inside a bounded exponent box.

\subsection{Exact counting from one local slope}

Let $G$ be a finite group, let $P_0\in G$ have order $d$, and suppose $Q_0=P_0^c$.  For integers $U,V\ge1$ consider
\[
 \psi_{U,V}(a,b)=P_0^aQ_0^b=P_0^{a+cb},
 \qquad 0\le a<U,\quad 0\le b<V.
\]

\begin{proposition}[Local-slope image bound]
\label{prop:slope-image}
Choose an integer representative $c$ of the local slope.  Then
\[
 \#\psi_{U,V}([0,U)\times[0,V))
 \le \min\{d,\ U+|c|(V-1)\}.
\]
If $0\le c<U$, the relation $(-c,1)$ produces at least
\[
 (U-c)(V-1)
\]
row collisions in the sense of Theorem~\ref{thm:short-relation}.
\end{proposition}

\begin{proof}
The integers $a+cb$ lie in an interval of length at most $U-1+|c|(V-1)$, so they assume at most $U+|c|(V-1)$ values before reduction modulo $d$, and never more than $d$ values after reduction.  If $0\le c<U$, then
\[
 P_0^{-c}Q_0=1.
\]
The two boxes $[0,U)\times[0,V)$ and its translate by $(-c,1)$ overlap in exactly $(U-c)(V-1)$ points.  Apply Theorem~\ref{thm:short-relation}.
\end{proof}

For the Alaoglu--Erd\H{o}s pair this says that a congruence
\[
 (M,A)\equiv(2^c,3^c)\pmod {p^t}
\]
with $0\le c<U$ contributes the factor
\[
 p^{t(U-c)(V-1)}
\]
to the row-grid determinant.  The same assertion can be read directly from a two-coordinate gcd.

\begin{definition}[Slope gcd]
For $c\in\Nzero$ put
\[
 H_c:=\gcd(|M-2^c|,|A-3^c|).
\]
For a nonintegral counterexample, $H_c$ is always a positive integer: $M=2^c$ and $A=3^c$ cannot occur simultaneously.
\end{definition}

\begin{corollary}[Weighted local-slope divisor]
\label{cor:slope-divisor}
Let $D$ be the determinant associated with the row box $[0,U)\times[0,V)$ and any equally large admissible column support.  For every $0\le c<U$,
\[
 H_c^{(U-c)(V-1)}\mid D.
\]
Consequently
\begin{equation}
 \mathfrak C_{U,V}
 :=\lcm_{0\le c<U} H_c^{(U-c)(V-1)}
 \quad\text{satisfies}\quad
 \mathfrak C_{U,V}\mid D.
 \label{eq:slope-lcm}
\end{equation}
\end{corollary}

\begin{proof}
If $p^t\mid H_c$, then $Q\equiv P^c\pmod {p^t}$.  Proposition~\ref{prop:slope-image} and Theorem~\ref{thm:short-relation} give
$v_p(D)\ge t(U-c)(V-1)$.  This holds prime by prime.  Taking the least common multiple is the exact way to combine possibly overlapping prime factors.
\end{proof}

\subsection{A fixed slope cannot recur indefinitely}

A key rigidity point is elementary but decisive.  Any local slopes capable of contributing at unbounded primes must move.

\begin{theorem}[Fixed-slope rigidity]
\label{thm:fixed-slope}
Fix $c\in\Z$.  If
\[
 \bar Q_p=\bar P_p^c
\]
for infinitely many good primes $p$, then $Q=P^c$ in $T(\Q)$.  In particular, under a nonintegral counterexample every fixed slope occurs at only finitely many good primes.
\end{theorem}

\begin{proof}
For $c\ge0$, every such prime divides both $M-2^c$ and $A-3^c$.  Infinitely many prime divisors force these two fixed integers to vanish, so $M=2^c$ and $A=3^c$.  For $c<0$, multiply the congruences by $2^{-c}$ and $3^{-c}$: every such prime divides both $2^{-c}M-1$ and $3^{-c}A-1$, and the same argument applies.  The conclusion $Q=P^c$ makes $\beta=c$.
\end{proof}

The proof also gives an exact weighted version.  For $c\ge0$,
\begin{equation}
 \sum_{\substack{p\ \mathrm{good}\\ \bar Q_p=\bar P_p^c}}\log p
 \le \log\rad H_c.
 \label{eq:fixed-slope-weight}
\end{equation}
Thus a successful finite-field argument cannot repeatedly reuse a bounded collection of slopes.  It must extract a large logarithmic weight from slopes $c_p$ that drift with $p$, or from high powers $p^t$ dividing the corresponding $H_{c_p}$.

\subsection{An intrinsic ceiling for the all-short-slope lcm}

The divisor~\eqref{eq:slope-lcm} is exact, but its maximum possible size can be bounded without any transcendence theory.  This produces a useful no-go theorem for the most naive short-slope strategy.

\begin{proposition}[Archimedean ceiling for the slope certificate]
\label{prop:slope-ceiling}
For fixed $M,A$ and $U,V\to\infty$,
\begin{align}
 \log \mathfrak C_{U,V}
 &\le (V-1)\sum_{c=0}^{U-1}(U-c)\log^+ H_c \notag\\
 &\le (V-1)\sum_{c=0}^{U-1}(U-c)\log^+|M-2^c| \notag\\
 &=\frac{\log2}{6}\,VU^3+O_M(VU^2).
 \label{eq:slope-ceiling}
\end{align}
Here $\log^+y=\max(0,\log y)$.
\end{proposition}

\begin{proof}
The logarithm of a least common multiple is bounded by the logarithm of the corresponding product, which gives the first line.  Since $H_c\le |M-2^c|$ unless the latter vanishes, and simultaneous vanishing is excluded, the second line follows after harmless treatment of finitely many small $c$.  Finally
\[
 \log|M-2^c|=c\log2+O_M(2^{-c})
\]
for large $c$, while
\[
 \sum_{c=0}^{U-1}(U-c)c=\frac{U(U-1)(U+1)}6.
\]
\end{proof}

Suppose the row box is aspect-optimized with area $m=UV$ and
$U^2\sim\beta m$, as in Section~\ref{sec:closure}.  Then Proposition~\ref{prop:slope-ceiling} becomes
\begin{equation}
 \log \mathfrak C_{U,V}
 \le\left(\frac{\beta\log2}{6}+o(1)\right)m^2.
 \label{eq:slope-ceiling-optimized}
\end{equation}
The optimized Hadamard budget below has leading constant
$4\sqrt{\beta\log2\log3}$.  Hence the \emph{maximum possible upper envelope} of this particular lcm certificate is smaller than that Hadamard constant whenever
\begin{equation}
 \beta<576\log_2 3=912.96\ldots.
 \label{eq:slope-threshold}
\end{equation}
This comparison does not upper-bound the true $p$-adic valuation of the determinant; independent relations and higher Smith layers can contribute more.  It does prove that merely listing all congruences $Q\equiv P^c$ with $0\le c<U$ and combining their individual gcds by an lcm cannot, by itself, close the present Hadamard estimate throughout the range~\eqref{eq:slope-threshold}.

\begin{remark}[Why this matters]
The finite computation in Section~\ref{sec:computation} finds exactly the behavior predicted by Theorem~\ref{thm:fixed-slope}.  For the genuine integral pair $Q=P^3$, the same slope appears at every good prime.  For an extremely close real near-miss, the few local slopes that occur are isolated and unrelated.  Archimedean closeness alone does not create a stable finite-field slope.
\end{remark}

\section{The finite contradiction criterion and its quantitative scale}
\label{sec:closure}

The preceding sections supply a lower bound for a positive integer determinant.  To turn it into a contradiction one needs an upper bound for the same determinant.  The deliberately coarse Hadamard estimate is enough to formulate an exact finite certificate and to expose the missing scale.

\subsection{An explicit archimedean budget}

Let
\[
 \Omega_{U,V}=\{0,\ldots,U-1\}\times\{0,\ldots,V-1\},
 \qquad
 \Sigma_{R,S}=\{0,\ldots,R-1\}\times\{0,\ldots,S-1\},
\]
and suppose
\[
 m:=UV=RS.
\]
Write $D_{U,V;R,S}$ for the determinant of the corresponding square evaluation matrix.  Under a hypothetical nonintegral solution, total positivity gives $D_{U,V;R,S}>0$.

Set
\begin{equation}
 T_{U,V}:=(U-1)+\beta(V-1),
 \qquad
 L_{R,S}:=(R-1)\log2+(S-1)\log3.
 \label{eq:TL}
\end{equation}
Every entry has the form $e^{t\lambda}$ with $0\le t\le T_{U,V}$ and $0\le\lambda\le L_{R,S}$.

\begin{theorem}[Hadamard budget]
\label{thm:hadamard}
One has
\begin{equation}
 0<\log D_{U,V;R,S}
 \le mT_{U,V}L_{R,S}+\frac m2\log m.
 \label{eq:hadamard-budget}
\end{equation}
\end{theorem}

\begin{proof}
Each row contains $m$ positive entries, all at most $e^{T_{U,V}L_{R,S}}$.  Its Euclidean norm is therefore at most
$\sqrt m\,e^{T_{U,V}L_{R,S}}$.  Hadamard's inequality gives
\[
 D_{U,V;R,S}
 \le\left(\sqrt m\,e^{T_{U,V}L_{R,S}}\right)^m.
\]
The determinant is positive by the exponential-kernel argument in Section~\ref{sec:grid}.
\end{proof}

\subsection{Exact orbit-collapse energy}

For every good prime power define
\[
 c^{\mathrm{row}}_{p,t}(U,V)
 :=\#\iota_{p^t}(\Omega_{U,V}),
 \qquad
 c^{\mathrm{col}}_{p,t}(R,S)
 :=\#\iota^*_{p^t}(\Sigma_{R,S}).
\]
The associated energy is
\begin{equation}
 \mathscr A_{U,V;R,S}
 :=\sum_{p\nmid6MA}\sum_{t\ge1}
 \left(m-\min\{c^{\mathrm{row}}_{p,t}(U,V),
                  c^{\mathrm{col}}_{p,t}(R,S)\}\right)\log p.
 \label{eq:orbit-energy}
\end{equation}
Only finitely many summands are nonzero for fixed boxes: once the reductions of all rows and columns are pairwise distinct modulo $p^t$, the corresponding term vanishes.

\begin{theorem}[Finite orbit-collapse certificate]
\label{thm:finite-certificate}
If, for some positive integers $U,V,R,S$ with $UV=RS=m$,
\begin{equation}
 \mathscr A_{U,V;R,S}
 >mT_{U,V}L_{R,S}+\frac m2\log m,
 \label{eq:finite-certificate}
\end{equation}
then no nonintegral Alaoglu--Erd\H{o}s solution can have output pair $(M,A)$.
\end{theorem}

\begin{proof}
The Smith tomography theorem gives
\[
 \log D_{U,V;R,S}\ge \mathscr A_{U,V;R,S},
\]
whereas Theorem~\ref{thm:hadamard} gives the opposite upper bound.  Strict inequality~\eqref{eq:finite-certificate} is impossible.
\end{proof}

This is a proof-producing criterion.  It involves only exact modular exponentiation, counting finite images, and a rigorously bounded real expression.  A finite list of primes and levels already gives a valid lower bound by truncating~\eqref{eq:orbit-energy}; no unproved distribution statement is needed to certify any successful instance.

\subsection{Optimizing the rectangle shapes}

The row and column boxes should not generally be square.  For fixed area $m$, the leading parts of $T_{U,V}$ and $L_{R,S}$ are minimized separately by balancing their two summands:
\begin{align}
 U&\sim\sqrt{\beta m},& V&\sim\sqrt{m/\beta},\label{eq:row-opt}\\
 R&\sim\sqrt{m\log_2 3},&
 S&\sim\sqrt{m\log_3 2}.\label{eq:col-opt}
\end{align}
Indeed, the arithmetic--geometric mean inequality gives
\begin{align}
 U+\beta V&\ge2\sqrt{\beta UV}=2\sqrt{\beta m},\notag\\
 R\log2+S\log3&\ge2\sqrt{RS\log2\log3}
 =2\sqrt{m\log2\log3}.
\end{align}
After rounding the four side lengths through a sequence of admissible factorizations of $m$,
\begin{equation}
 mT_{U,V}L_{R,S}
 =\left(4\sqrt{\beta\log2\log3}+o(1)\right)m^2.
 \label{eq:optimized-hadamard}
\end{equation}
The optimization reduces the large-$\beta$ dependence from linear to square-root, but the scale remains $m^2$.

\subsection{Why cyclic image size alone is not enough}

Suppose, unrealistically favorably, that for every good prime power the full row image were cyclic.  Its size would then be at most $\varphi(p^t)$.  Replacing the exact orbit count by this bound yields the formal lower certificate
\begin{equation}
 \mathscr C(m)
 :=\sum_p\sum_{t\ge1}(m-\varphi(p^t))_+\log p.
 \label{eq:cyclic-cardinality-certificate}
\end{equation}

\begin{proposition}[Cyclic-cardinality ceiling]
\label{prop:cyclic-ceiling}
As $m\to\infty$,
\begin{equation}
 \mathscr C(m)=\frac12m^2+o(m^2).
 \label{eq:cyclic-ceiling}
\end{equation}
\end{proposition}

\begin{proof}
The level $t=1$ contributes
\[
 \sum_{p\le m+1}(m-p+1)\log p
 =m\vartheta(m)-\sum_{p\le m}p\log p+o(m^2),
\]
where $\vartheta$ is Chebyshev's function.  Partial summation and the prime number theorem give
$\vartheta(m)\sim m$ and $\sum_{p\le m}p\log p\sim m^2/2$, hence the main term $m^2/2$.  For $t\ge2$, the condition $\varphi(p^t)<m$ forces $p\ll m^{1/t}$; summing the trivial bound $m\log p$ gives $O(m^{3/2})$ at $t=2$ and smaller orders thereafter.
\end{proof}

Compare the constant $1/2$ with the optimized Hadamard constant
$4\sqrt{\beta\log2\log3}$.  Even for $\beta=1$, the latter exceeds $3.49$.  Thus \emph{cyclic reduction image at every finite level} is quantitatively insufficient under the present upper bound.  The determinant must see more than local rank one: it needs short slopes, repeated high-power congruence, multiple independent relations, or a significantly sharper archimedean estimate.

\subsection{A quantitative moving-slope target}

The exact short-relation theorem turns the preceding diagnosis into a clean sufficient condition.  Fix $0<\eta<1$ and let $\mathcal P_m$ be a set of good primes for which there exists an integer $c_p$ satisfying
\[
 |c_p|\le\eta U,
 \qquad
 \bar Q_p=\bar P_p^{c_p}.
\]
Choose the sign convention so that the overlap of the row box with its translate by $(-c_p,1)$ is at least $(1-\eta)U(V-1)$.  Then
\begin{equation}
 \mathscr A_{U,V;R,S}
 \ge (1-\eta)U(V-1)
     \sum_{p\in\mathcal P_m}\log p.
 \label{eq:moving-slope-lower}
\end{equation}
Since $U(V-1)=m+O(U)$, an aspect-optimized sequence would contradict the hypothetical pair if
\begin{equation}
 \sum_{p\in\mathcal P_m}\log p
 >\left(\frac{4\sqrt{\beta\log2\log3}}{1-\eta}+o(1)\right)m.
 \label{eq:moving-slope-target}
\end{equation}
A $p^t$ congruence receives the weight $t\log p$, so the same criterion admits depth in place of more primes.

Condition~\eqref{eq:moving-slope-target} is the first genuinely quantitative target of this report.  It asks for logarithmic prime weight of order $m$, not density among all primes, and it explicitly allows the slopes to drift as Theorem~\ref{thm:fixed-slope} requires.  What remains unknown is a theorem transferring the singular real logarithm matrix into this much finite-place short-slope mass.

\section{Local--global support theorems specialized to the conjecture}
\label{sec:support}

The finite-field framework naturally meets the classical support problem for algebraic groups.  The relevant results are due to Corrales--Rodrig\'a\~nez and Schoof, and in the general product-of-tori-and-abelian-varieties form to Perucca~\cite{CorralesSchoof,PeruccaSupport,PeruccaDetect,PeruccaOrders}.  Their qualitative conclusions are strong enough to give exact conditional proofs of Alaoglu--Erd\H{o}s under almost-everywhere local hypotheses.  They do not supply the quantitative short relations required by~\eqref{eq:moving-slope-target}.

\subsection{Cyclic subgroup detection}

We use the following imported theorem in its torus specialization.

\begin{theorem}[Perucca, cyclic reduction detection]
\label{thm:perucca-cyclic}
Let $G$ be a product of an abelian variety and a torus over a number field, let $R\in G(K)$, and let $\Lambda\subset G(K)$ be a free cyclic subgroup.  If the reduction of $R$ belongs to the reduction of $\Lambda$ for all but finitely many primes of good reduction, then $R\in\Lambda$.
\end{theorem}

Apply this with $G=T=(\Gm)^2$, $R=Q$, and $\Lambda=\langle P\rangle$.

\begin{corollary}[Almost-everywhere local cyclicity proves Alaoglu--Erd\H{o}s]
\label{cor:local-cyclic-AE}
If
\[
 \bar Q_p\in\langle\bar P_p\rangle
\]
for all but finitely many good primes $p$, then $\beta\in\Z$.
\end{corollary}

\begin{proof}
Theorem~\ref{thm:perucca-cyclic} gives $Q=P^n$ for some $n\in\Z$.  Positivity and $M,A>1$ give $n>0$, and then $M=2^n$, $A=3^n$, so $\beta=n$.
\end{proof}

The contrapositive is unconditional and useful:
\begin{equation}
 \beta\notin\Z
 \quad\Longrightarrow\quad
 \bar Q_p\notin\langle\bar P_p\rangle
 \text{ for infinitely many good primes }p.
 \label{eq:noncyclic-witnesses}
\end{equation}
It says that a counterexample cannot be locally rank one almost everywhere.  Qualitative support theory therefore creates infinitely many \emph{anti-collapse witnesses}.  A proof via determinant collapse must coexist with these witnesses and concentrate on a sufficiently large complementary set.

\subsection{Order support already suffices}

A second support theorem uses only the orders of the reductions.  In a torus, endomorphisms are integral matrices.

\begin{theorem}[Perucca, order support; torus form]
\label{thm:perucca-order}
Let $P,Q\in T(\Q)$.  If
\[
 \ord(\bar Q_p)\mid\ord(\bar P_p)
\]
for all but finitely many good primes, then there exist a nonzero integer $c$ and an endomorphism
$\phi\in\operatorname{End}_{\Q}(T)\cong M_2(\Z)$ such that
\[
 \phi(P)=Q^c.
\]
\end{theorem}

\begin{corollary}[Almost-everywhere order divisibility proves the conjecture]
\label{cor:order-AE}
Under the hypothetical output relation $M=2^\beta$, $A=3^\beta$, if
\[
 \ord(\bar Q_p)\mid\ord(\bar P_p)
\]
for all but finitely many good primes, then $\beta\in\Z$.
\end{corollary}

\begin{proof}
Write
\[
 \phi=\begin{pmatrix}u&v\\ r&s\end{pmatrix}.
\]
The equation $\phi(P)=Q^c$ says
\[
 M^c=2^u3^v,
 \qquad
 A^c=2^r3^s.
\]
Every prime divisor of $M$ or $A$ is therefore $2$ or $3$, so write
\[
 M=2^{u_0}3^{v_0},
 \qquad A=2^{r_0}3^{s_0}
\]
with nonnegative exponents.  Put $\theta=\log_2 3$.  The equality $2^\beta=M$ gives
$\beta=u_0+v_0\theta$, while $3^\beta=A$ gives, after taking logarithms base $2$,
\[
 \beta\theta=r_0+s_0\theta.
\]
Eliminating $\beta$ yields
\[
 v_0\theta^2+(u_0-s_0)\theta-r_0=0.
\]
The number $\theta$ is transcendental: it is irrational by unique factorization, and if it were algebraic then the Gelfond--Schneider theorem applied to $2^\theta=3$ would be contradicted.  Hence every coefficient in the displayed polynomial is zero.  Thus $v_0=r_0=0$ and $u_0=s_0$, so $M=2^{u_0}$, $A=3^{u_0}$ and $\beta=u_0$.
\end{proof}

Again the contrapositive supplies an infinite local signature:
\begin{equation}
 \beta\notin\Z
 \quad\Longrightarrow\quad
 \ord(\bar Q_p)\nmid\ord(\bar P_p)
 \text{ for infinitely many good }p.
 \label{eq:order-witnesses}
\end{equation}
Perucca also proves variants that compare the $\ell$-adic valuations of the two orders for a fixed auxiliary prime $\ell$~\cite{PeruccaSupport}.  Those variants are particularly relevant to the shape of the congruence lattice.

\subsection{Order imbalance forces long coefficients}

\begin{lemma}[Order-imbalance obstruction]
\label{lem:order-imbalance}
Let $G$ be a finite abelian group, $P_0,Q_0\in G$, and suppose
\[
 P_0^aQ_0^b=1.
\]
For every prime $\ell$,
\begin{equation}
 v_\ell(b)\ge
 \bigl(v_\ell(\ord Q_0)-v_\ell(\ord P_0)\bigr)_+.
 \label{eq:order-imbalance}
\end{equation}
The symmetric assertion holds with $a$ and $b$ interchanged.
\end{lemma}

\begin{proof}
The relation gives $Q_0^b=P_0^{-a}$, whose order divides $\ord P_0$.  Since
\[
 v_\ell(\ord(Q_0^b))
 =\max\{v_\ell(\ord Q_0)-v_\ell(b),0\},
\]
comparison with $v_\ell(\ord P_0)$ gives~\eqref{eq:order-imbalance}.
\end{proof}

This lemma shows why support-theorem witness primes need not help the determinant.  If the $\ell$-part of $\ord(\bar Q_p)$ exceeds that of $\ord(\bar P_p)$ by $j$, then every local relation with nonzero $Q$-coefficient has that coefficient divisible by $\ell^j$.  For a box with $V\le\ell^j$, no such relation is short in the vertical direction.  The same primes that certify global nondependence can therefore make the reduction lattice \emph{less} compressible.

\begin{resultbox}
Qualitative local--global theory reaches an exact logical endpoint:
\begin{itemize}
\item local cyclicity almost everywhere would prove $\beta\in\Z$;
\item local order domination almost everywhere would prove $\beta\in\Z$;
\item a counterexample must have infinitely many noncyclic and order-imbalance witness primes;
\item none of these statements controls the logarithmic mass of primes carrying \emph{short} relations.
\end{itemize}
The missing theorem is quantitative and geometric in the two-dimensional relation lattice, not merely qualitative membership in a subgroup.
\end{resultbox}

\section{Fixed-ray common divisors and why moving directions are essential}
\label{sec:fixed-rays}

The weighted gcd divisors of Section~\ref{sec:gcd} might suggest taking one promising relation direction and scaling it.  The theorem of Bugeaud, Corvaja, and Zannier rules out that architecture at exponential scale~\cite{BCZ2003}.

\subsection{The row rays}

Fix $(h,k)\in\Nzero^2\setminus\{(0,0)\}$ and put
\[
 U_0=2^hM^k,
 \qquad V_0=3^hA^k.
\]
Under a hypothetical counterexample,
\[
 \log U_0=(h+k\beta)\log2,
 \qquad
 \log V_0=(h+k\beta)\log3.
\]
The integers $U_0$ and $V_0$ are multiplicatively independent.  Indeed, a relation $U_0^r=V_0^s$ would imply $r\log2=s\log3$ because $h+k\beta>0$, which is impossible unless $r=s=0$.

The Bugeaud--Corvaja--Zannier theorem therefore gives, for every $\varepsilon>0$ and all sufficiently large $n$,
\begin{equation}
 \log\gcd(U_0^n-1,V_0^n-1)<\varepsilon n.
 \label{eq:BCZ-row}
\end{equation}
But this gcd is exactly
\[
 G^{\mathrm{row}}_{nh,nk}.
\]
Thus no fixed row direction carries a common divisor of positive exponential rate.

\subsection{The column rays}

Fix $(r,s)\in\Nzero^2\setminus\{(0,0)\}$ and put
\[
 W_0=2^r3^s,
 \qquad Z_0=M^rA^s=W_0^\beta.
\]
These are multiplicatively independent: if $W_0^u=Z_0^v$, then
$W_0^{u-v\beta}=1$, so $u=v\beta$; irrationality of $\beta$ forces $u=v=0$.  Therefore
\begin{equation}
 \log\gcd(W_0^n-1,Z_0^n-1)<\varepsilon n
 \label{eq:BCZ-column}
\end{equation}
for large $n$, and the gcd is $G^{\mathrm{col}}_{nr,ns}$.

\subsection{A finite-family no-go theorem}

\begin{proposition}[Finitely many fixed rays are lower order]
\label{prop:fixed-rays}
Fix finitely many nonzero row and column directions.  Let the determinant boxes have side lengths $O(N)$ in each coordinate, and combine the weighted gcd divisors from \emph{all} integer multiples of those directions that fit in the boxes.  The logarithm of the resulting lcm is $o(N^4)$.  In particular it is $o(m^2)$ when the matrix size is $m\asymp N^2$.
\end{proposition}

\begin{proof}
For one fixed direction, a multiple $n$ that remains inside the box satisfies $n=O(N)$, and its overlap weight is $O(N^2)$.  By~\eqref{eq:BCZ-row} or~\eqref{eq:BCZ-column}, if $g_n$ denotes the associated gcd, then $\log g_n=o(n)$.  Hence
\[
 \sum_{n\le C N}\log g_n=o(N^2).
\]
The logarithm of the lcm of the weighted divisors is at most the logarithm of their product, and therefore at most
$O(N^2)\sum_{n\le C N}\log g_n=o(N^4)$.  A fixed finite sum over directions remains $o(N^4)$.
\end{proof}

The conclusion is stronger than the statement that one particular gcd is small.  It says that any closing common-divisor mechanism must use a number of relation directions growing with the determinant, and those directions must move through the lattice.  This matches independently the fixed-slope rigidity of Theorem~\ref{thm:fixed-slope} and the conditional target~\eqref{eq:moving-slope-target}.

\begin{remark}[Uniformity is the open issue]
The Bugeaud--Corvaja--Zannier bound is asymptotic for fixed $U_0,V_0$.  The bases in the moving-direction problem themselves grow with the box.  A uniform theorem over two-parameter families, with information about the \emph{union} of prime divisors rather than a single gcd, would be a genuinely new input.  Existing $S$-unit and Subspace-Theorem technology does not provide the needed lower or upper control at this moving scale.
\end{remark}

\section{Exact finite-field reconnaissance}
\label{sec:computation}

Theorems~\ref{thm:smith-tomography} and~\ref{thm:finite-certificate} are exact, but they leave open how exceptional the orbit collapse is for actual integer pairs.  I wrote a dependency-free Python scout to measure the one-level row image
\[
 \{(2^aM^b,3^aA^b)\bmod p:0\le a,b<N\}
\]
for every good prime $p\le5000$.  The purpose was diagnostic, not evidentiary: the calculation does not assume the conjecture, and no numerical observation is promoted to a theorem.

\subsection{Protocol}

The run used $N=24$, hence $m=N^2=576$.  For each pair $(M,A)$ and each good prime, the program computed:
\begin{enumerate}[label=\textup{(\roman*)}]
\item the exact row-orbit image size $c_p$ and defect $m-c_p$;
\item the orders of $P=(2,3)$ and $Q=(M,A)$ in $(\F_p^\times)^2$;
\item whether $Q\in\langle P\rangle$, and if so the least nonnegative slope;
\item a primitive-root discrete-log matrix $B_p$;
\item the exact image size $|(\Z/(p-1)\Z)^2B_p|$ from the determinantal formula~\eqref{eq:image-formula};
\item whether $p-1\mid\det B_p$, the rank-one-scale condition of Corollary~\ref{cor:rank-one-scale}.
\end{enumerate}
The direct box enumeration and the Smith/determinantal image formula agreed for every row of the data set.

The four test pairs were:
\begin{align*}
 &(8,27), &&\text{the exact integral pair }(2^3,3^3);\\
 &(5,13), &&\text{a generic integer control};\\
 &(3779,467744), &&\text{an exceptionally close real-logarithm near-miss};\\
 &(9,2), &&\text{a deliberately asymmetric control}.
\end{align*}
For the near-miss,
\begin{align}
 \log_2(3779)&=11.883026394927862\ldots,\notag\\
 \log_3(467744)&=11.883026393570518\ldots,
 \label{eq:near-miss}
\end{align}
whose discrepancy is only
$1.3573444583858671\times10^{-9}$.

\subsection{Results}

Let ``large defect'' mean defect at least $m/4=144$.  The quantity in the penultimate column below is the one-level row sum
\[
 \mathscr E_{\le5000}^{\mathrm{row}}
 :=\sum_{\substack{p\le5000\\p\nmid6MA}}(m-c_p)\log p.
\]

\begin{table}[htbp]
\centering
\small
\setlength{\tabcolsep}{4.5pt}
\begin{tabular}{@{}lrrrrrrrr@{}}
\toprule
pair $(M,A)$ & good $p$ & collapse & large & local & rank-one & $\det=0$ & max & $\mathscr E$\\
&& primes & defect & cyclic & scale & mod $p-1$ & defect & \\
\midrule
$(8,27)$ & 667 & 667 & 667 & 667 & 667 & 667 & 572 & $2{,}376{,}056.63$\\
$(5,13)$ & 665 & 13 & 11 & 1 & 2 & 1 & 540 & $12{,}425.85$\\
$(3779,467744)$ & 664 & 15 & 14 & 2 & 2 & 1 & 572 & $15{,}007.14$\\
$(9,2)$ & 667 & 14 & 11 & 0 & 0 & 0 & 552 & $13{,}122.57$\\
\bottomrule
\end{tabular}
\caption{Exact one-level row-orbit statistics for $N=24$ and good primes $p\le5000$.  ``Collapse'' means $c_p<m$.  ``Local cyclic'' means $\bar Q_p\in\langle\bar P_p\rangle$.}
\label{tab:scout}
\end{table}

The integral pair is separated perfectly: every good prime is locally cyclic with the same slope $3$, every discrete-log determinant vanishes modulo $p-1$, and every box collapses.  This is the finite-field shadow of the global identity $Q=P^3$; the associated integer grid itself has repeated rows, so this control case is not an admissible positive determinant.

The three nonintegral controls look qualitatively alike.  Almost every prime has the full $24\times24$ box image.  Large collapses are rare, and local cyclicity is rarer still.  In particular, the $10^{-9}$ agreement in~\eqref{eq:near-miss} does not produce a corresponding broad modular signature.  Among the exceptional primes:
\begin{itemize}
\item for $(5,13)$, $p=37$ has local slope $30$, image size $36$, and defect $540$; $p=47$ has rank-one-scale image size $46$ but no local slope;
\item for the near-miss, $p=5$ has local slope $2$, image size $4$, and defect $572$; $p=379$ has local slope $362$, image size $189$, and defect $387$.
\end{itemize}
The second near-miss slope is large compared with the side length $24$, so it gives no direct short translation inside the test box.  This illustrates the difference between local cyclicity and useful local cyclicity.

\subsection{What the experiment establishes and what it does not}

The run verifies five implementation-level facts relevant to a future proof search:
\begin{enumerate}[label=\textup{(\arabic*)}]
\item the direct orbit count agrees exactly with the discrete-log Smith formula;
\item transpose duality predicts identical full-subgroup cardinalities, even though truncated row and column boxes may behave differently;
\item exact integral dependence creates a uniform, prime-by-prime signature;
\item extremely accurate real-logarithm matching does not by itself create that signature;
\item the informative exceptional primes often carry large slopes, confirming the need to record relation length rather than only image rank.
\end{enumerate}
It does \emph{not} test prime powers $p^t$, does not enumerate the column box, and does not approach the asymptotic regime required by~\eqref{eq:moving-slope-target}.  It also cannot rule out a sparse but eventually decisive family of primes beyond $5000$.

\subsection{Reproducibility record}

The run produces \texttt{finite\_field\_scout.csv} and
\texttt{finite\_field\_scout\_summary.json}. The complete Python source is embedded in
Appendix~\ref{app:code}; it uses only the Python standard library.

\section{A formalization-ready decomposition}
\label{sec:formalization}

The finite-field route has an advantage over several analytic approaches in the source report: its central implications can be decomposed into small algebraic statements with explicit certificates.  This section records a Lean-oriented architecture.  It is a design, not a claim that the formalization has already been completed.

\subsection{Core objects}

A formal development can avoid transcendental matrix entries until the final positivity lemma.  Fix natural numbers $M,A,U,V,R,S$ and define:
\begin{enumerate}[label=\textup{(F\arabic*)}]
\item finite index types for the row and column boxes;
\item the integer entry function
\[
 E(a,b,r,s)=(2^aM^b)^r(3^aA^b)^s;
\]
\item the square integer matrix and its determinant $D$;
\item for each modulus $q$, the reduction maps $\iota_q$ and $\iota_q^*$;
\item finite sets representing their box images.
\end{enumerate}
All exponentiation in this layer is natural-number exponentiation.  The only assumptions needed for modular inverses are explicit coprimality hypotheses.

\subsection{The easiest certified divisor: one short relation}

Theorem~\ref{thm:short-relation} admits a direct formal proof that does not require a general Smith-normal-form library:
\begin{enumerate}[label=\textup{(L\arabic*)}]
\item choose a spanning forest of the directed translation graph on the box;
\item order vertices so that every non-root follows its parent;
\item replace each non-root row by row minus parent row;
\item prove that every entry in each replaced row is divisible by $p^t$;
\item factor $p^t$ from each such row and use determinant invariance under row addition.
\end{enumerate}
The number of non-root vertices is exactly the overlap count.  This yields a compact checker whose input is only $(p,t,h,k)$ plus two modular exponentiation equalities.

\subsection{Tomography without choosing a Smith normal form}

The full layer theorem can be formalized in either of two ways.

\paragraph{Smith route.}
Use the Smith normal form over $\Z$ and show that the image of the reduced diagonal matrix over $\Z/p^t\Z$ has minimal generator number
$\#\{i:v_p(s_i)<t\}$.  Bound that generator number by the number of distinct rows or columns.

\paragraph{Determinantal-ideal route.}
For each $j$, show that if the reduced row module is generated by at most $c$ elements modulo $p^t$, then every $(c+1)\times(c+1)$ minor is divisible by $p^t$.  Apply the relation between determinantal ideals and partial products of invariant factors, then sum over $t$.  This route may align better with existing commutative-algebra infrastructure.

The short-relation certificate should be formalized first because it already captures the moving-slope target and has a smaller dependency surface.

\subsection{The positivity boundary}

To use a nonzero determinant, the formal development needs the fact that
\[
 \det(e^{x_i y_j})_{1\le i,j\le m}>0
\]
for strictly increasing real sequences $x_i$ and $y_j$.  Several proof paths are available:
\begin{enumerate}[label=\textup{(P\arabic*)}]
\item invoke strict total positivity of the exponential kernel;
\item prove the generalized Vandermonde determinant is positive by repeated integral factorization;
\item use the Chebyshev-system criterion for the functions $e^{y_jx}$.
\end{enumerate}
For this application, strict ordering follows from irrationality of $\beta$ and irrationality of $\log_3 2$.  Irrationality of $\beta$ is elementary from the integrality assumptions: a rational exponent producing an integer power of $2$ must be integral.

\subsection{A machine-checkable terminal certificate}

A finite successful certificate could have the following serialized form:
\begin{verbatim}
(M, A, U, V, R, S,
 [(p, t, row_image_partition, column_image_partition), ...],
 rational_upper_bound)
\end{verbatim}
The checker would:
\begin{enumerate}[label=\textup{(C\arabic*)}]
\item verify primality and good reduction;
\item recompute every modular row and column image;
\item sum the exact Smith-layer lower bound;
\item verify an outward-rounded rational upper bound for
$mT_{U,V}L_{R,S}+\tfrac m2\log m$;
\item compare the two rational numbers.
\end{enumerate}
No trust in the search program would be required.  The search can be heuristic; only the small certificate checker needs formal verification.

\section{Research targets produced by the finite-field route}
\label{sec:targets}

The route does not yet prove the conjecture.  It does, however, replace a diffuse search by several sharply falsifiable statements.  Each target below comes with a precise way it would close, or fail to close, the determinant argument.

\subsection{Target I: moving-slope logarithmic mass}

Let $U_m,V_m$ be aspect-optimized row sides with $U_mV_m=m$, and for $0<\eta<1$ define
\[
 \mathcal P_m(\eta)
 :=\left\{p\nmid6MA:\exists c\in\Z,
 |c|\le\eta U_m,
 Q\equiv P^c\pmod p\right\}.
\]
A sufficient theorem is
\begin{equation}
 \limsup_{m\to\infty}
 \frac1m\sum_{p\in\mathcal P_m(\eta)}\log p
 >\frac{4\sqrt{\beta\log2\log3}}{1-\eta}.
 \label{eq:target-I}
\end{equation}
By~\eqref{eq:moving-slope-lower}, this would contradict a nonintegral pair.  Theorem~\ref{thm:fixed-slope} shows that any proof must exploit a growing collection of slopes.  A direct attack could study the prime divisors of the moving family
\[
 \gcd(M-2^c,A-3^c),\qquad |c|\le\eta U_m,
\]
with multiplicity and with overlap between different $c$ handled by the lcm rather than the product.

A falsifier would be a hypothetical pair for which the left side of~\eqref{eq:target-I} remains below the displayed threshold for every $\eta<1$.  Even such a negative result would be valuable because it would rule out the most direct one-relation determinant closure.

\subsection{Target II: $p$-adic depth without short level-one slopes}

Level-one local slopes may be sparse, but a prime can contribute through several Smith layers.  Define the best short-slope depth
\[
 d_p(U):=
 \max_{|c|<U}\min\{v_p(M-2^c),v_p(A-3^c)\},
\]
with the negative-$c$ expressions cleared of denominators.  The short-relation theorem gives roughly
\[
 (U-|c|)(V-1)d_p(U)\log p
\]
from the maximizing slope.  A depth version of Target I is therefore
\begin{equation}
 \sum_p d_p(U_m)\,(1-|c_p|/U_m)_+\log p
 >\left(4\sqrt{\beta\log2\log3}+o(1)\right)m.
 \label{eq:target-II}
\end{equation}
The challenge is simultaneous lifting in two coordinates.  Ordinary Hensel lifting of a discrete logarithm in one coordinate does not force the same lifted slope in the other.

\subsection{Target III: two independent short relations}

One local relation collapses a box along one direction.  If $L^{\mathrm{row}}_{p^t}$ contains two independent vectors $u_1,u_2$ whose fundamental parallelogram has small area relative to $UV$, then the orbit count can be much smaller than any one-overlap estimate.  Geometry of numbers suggests studying the successive minima
\[
 \lambda_1(L^{\mathrm{row}}_{p^t}),
 \qquad
 \lambda_2(L^{\mathrm{row}}_{p^t})
\]
with respect to the rectangle norm
$\|(h,k)\|_{U,V}=\max(|h|/U,|k|/V)$.  The exact theorem needed is a lower bound on
\[
 \sum_{p,t}
 \left(m-\#\bigl(\Omega_{U,V}/L^{\mathrm{row}}_{p^t}\bigr)\right)\log p
\]
that exploits both minima, not only the lattice index.  Discrete-log determinants control the index; the unresolved information is the shape.

\subsection{Target IV: sharpen the archimedean determinant}

Hadamard discards all cancellation and the precise total-positive structure.  For increasing $x_i,y_j$ there are exact factorizations
\[
 \det(e^{x_i y_j})
 =e^{\sum_i x_i y_1}
  \prod_{i<j}(x_j-x_i)
  \prod_{i<j}(y_j-y_i)
  \cdot \mathcal I(x,y),
\]
where $\mathcal I$ is a positive iterated-integral factor.  A uniform upper bound that saves a fixed proportion of the leading constant in~\eqref{eq:optimized-hadamard} would immediately lower the required local mass.  To be decisive against the cyclic-cardinality ceiling alone, the leading constant would have to fall below $1/2$, an enormous improvement; more plausibly, a moderate saving would combine with short-relation divisors.

\subsection{Target V: transfer singularity of logarithms to finite places}

The defining real relation is the rank-one identity
\[
 \det\begin{pmatrix}\log2&\log M\\\log3&\log A\end{pmatrix}=0.
\]
At a prime with chosen primitive root, the finite analogue is
\[
 \det B_p\equiv0\pmod{p-1}.
\]
A theorem forcing this congruence for a positive logarithmic mass of primes would already yield many rank-one-scale images, but Proposition~\ref{prop:cyclic-ceiling} shows that image size alone is not enough.  The transfer theorem must be refined to produce small kernel vectors or substantial $p$-adic depth.  This is the conceptual bottleneck:
\begin{quote}
convert an exact singularity in the real logarithm matrix into quantitative shape information for the kernels of its finite logarithm matrices.
\end{quote}
No known support theorem provides such a transfer.

\section{Conclusions}
\label{sec:conclusion}

The investigation did not produce a complete proof of the Alaoglu--Erd\H{o}s conjecture.  It produced a new finite-place framework with exact theorems, exact computational diagnostics, and several negative results that materially narrow the remaining search.

The principal advances are:
\begin{enumerate}[label=\textup{(\arabic*)}]
\item \emph{Pair-specific Smith tomography.}  The exact reductions of $(2,3)$ and $(M,A)$ modulo every $p^t$ bound every Smith layer of the integer grid determinant; this is stronger and more tailored than a universal polynomial-function divisor.
\item \emph{Short-relation amplification.}  A local relation $(h,k)$ contributes its geometric overlap multiplicity, yielding explicit weighted row and column gcd divisors.
\item \emph{Finite closure.}  The orbit-collapse energy~\eqref{eq:orbit-energy} gives a completely finite contradiction certificate against the Hadamard budget~\eqref{eq:hadamard-budget}.
\item \emph{Sharp route diagnostics.}  Cyclic image size at every prime-power level contributes only $(1/2+o(1))m^2$ under the present estimate; every fixed slope and every finite family of fixed relation rays are also insufficient at the leading $m^2$ scale.
\item \emph{Local--global boundary.}  Existing support theorems would prove integrality from almost-everywhere local cyclicity or order domination, but a counterexample must have infinitely many anti-collapse witness primes.  What is missing is quantitative control of the complementary short-relation mass.
\item \emph{Reproducible reconnaissance.}  Exact finite-field computations sharply distinguish a true integral-power pair from generic pairs and from an archimedean near-miss, while confirming the discrete-log image formula on every tested prime.
\end{enumerate}

The strongest concrete next statement is~\eqref{eq:target-I}, or its depth and two-relation variants.  A proof of one of these quantitative targets, together with Theorem~\ref{thm:finite-certificate}, would finish the conjecture.  Conversely, a systematic failure would establish that the determinant route requires a fundamentally sharper archimedean ingredient.  The problem has therefore been reduced, within this approach, to a precise local-shape question rather than another equivalent reformulation of four exponentials.

\appendix

\section{Expanded proof of the Smith-layer identity}
\label{app:smith}

For completeness, this appendix records the module-theoretic identity used in Theorem~\ref{thm:smith-tomography}.  Let $E\in M_m(\Z)$ be nonsingular and write its Smith form over $\Z_p$ as
\[
 UE V=\diag(p^{\nu_1},\ldots,p^{\nu_m}),
 \qquad \nu_1\le\cdots\le\nu_m,
\]
with $U,V\in\operatorname{GL}_m(\Z_p)$.  Reduction modulo $p^t$ preserves the isomorphism class of the row and column modules.  The row module of the diagonal matrix over $R_t=\Z/p^t\Z$ is
\[
 \bigoplus_{\nu_i<t}p^{\nu_i}R_t.
\]
Each summand is isomorphic to $R_t/p^{t-\nu_i}R_t$ and contributes one dimension to
$M/pM$ over $\F_p$.  Nakayama's lemma therefore gives
\[
 \mu_{R_t}(\operatorname{row}(E\bmod p^t))
 =\#\{i:\nu_i<t\}=r_{p,t},
\]
where $\mu$ denotes minimal generator number.  If only $c$ distinct rows occur, those rows generate the module, so $r_{p,t}\le c$.  The same argument applies to columns.

Finally, for every nonnegative integer $\nu$,
\[
 \nu=\sum_{t\ge1}\mathbf1_{\nu\ge t}.
\]
Summing over the Smith valuations gives
\[
 v_p(\det E)
 =\sum_i\nu_i
 =\sum_{t\ge1}\#\{i:\nu_i\ge t\}
 =\sum_{t\ge1}(m-r_{p,t}).
\]
This proves the layer-cake identity and explains why counting image classes at \emph{all} prime-power depths is exactly matched to determinant valuation.

\section{A direct elimination proof of the short-relation divisor}
\label{app:elimination}

Let $\Omega\subset\Z^2$ be finite and let $u\ne0$ satisfy
$P^{u_1}Q^{u_2}\equiv1\pmod{p^t}$.  Form a graph on $\Omega$ by joining $\omega$ and $\omega+u$ whenever both belong to $\Omega$.  Because translation by a nonzero vector has no finite cycles in $\Z^2$, this graph is a forest.  Choose one root in each component and orient every edge away from the root.

Order the rows so that parents precede children.  For every non-root vertex $\omega$, replace the row indexed by $\omega$ by
\[
 \operatorname{row}(\omega)-\operatorname{row}(\operatorname{parent}(\omega)).
\]
The operation preserves the determinant.  The two row parameters differ by $u$, so every entry in the new row is divisible by $p^t$.  Factoring $p^t$ from every non-root row gives
\[
 p^{t(|\Omega|-\#\text{components})}\mid D.
\]
A forest has $|\Omega|-\#\text{components}$ edges.  For a rectangle and a single translation $u=(h,k)$, that edge count is the exact overlap
$(U-|h|)_+(V-|k|)_+$.  This proof is especially convenient for formal verification because it avoids any appeal to Smith normal form.

\section{Exact reconnaissance code}
\label{app:code}

The following is the complete standard-library Python program used for Table~\ref{tab:scout}.  It writes both the row-level CSV file and the JSON summary.

\begin{lstlisting}[language=Python,basicstyle=\ttfamily\tiny,breaklines=true,breakatwhitespace=false,columns=fullflexible,keepspaces=true]
#!/usr/bin/env python3
"""Finite-field orbit scout for the Alaoglu--Erdos grid.

For each good prime p, the program computes:
  * the discrete-log matrix B_p for P=(2,3), Q=(M,A),
  * the full orbit size |im B_p| in (Z/(p-1)Z)^2,
  * exact row/column image counts for the N x N exponent box,
  * whether Q mod p lies in the cyclic subgroup generated by P mod p.

All arithmetic is exact integer/modular arithmetic.  This is exploratory evidence,
not a proof of the conjecture.
"""
from __future__ import annotations

import argparse
import csv
import json
import math
from pathlib import Path
from typing import Dict, Iterable, List, Sequence, Tuple

Pair = Tuple[int, int, str]


def primes_upto(limit: int) -> List[int]:
    if limit < 2:
        return []
    sieve = bytearray(b"\x01") * (limit + 1)
    sieve[0:2] = b"\x00\x00"
    stop = int(limit**0.5)
    for q in range(2, stop + 1):
        if sieve[q]:
            start = q * q
            sieve[start : limit + 1 : q] = b"\x00" * (((limit - start) // q) + 1)
    return [q for q in range(2, limit + 1) if sieve[q]]


def prime_factors(n: int) -> List[int]:
    ans: List[int] = []
    d = 2
    while d * d <= n:
        if n % d == 0:
            ans.append(d)
            while n % d == 0:
                n //= d
        d = 3 if d == 2 else d + 2
    if n > 1:
        ans.append(n)
    return ans


def primitive_root_prime(p: int) -> int:
    if p == 2:
        return 1
    n = p - 1
    factors = prime_factors(n)
    for g in range(2, p):
        if all(pow(g, n // ell, p) != 1 for ell in factors):
            return g
    raise RuntimeError(f"No primitive root found modulo {p}")


def discrete_log_table(p: int, g: int) -> Dict[int, int]:
    table: Dict[int, int] = {}
    x = 1
    for k in range(p - 1):
        table[x] = k
        x = (x * g) % p
    if len(table) != p - 1:
        raise RuntimeError(f"{g} is not primitive modulo {p}")
    return table


def image_size_mod_n(matrix: Tuple[int, int, int, int], n: int) -> int:
    """Cardinality of the image of a 2x2 integer matrix on (Z/nZ)^2.

    If B=(b_ij), the cokernel has order equal to the gcd of the 2x2 minors
    of [n I_2 | B].
    """
    b11, b12, b21, b22 = matrix
    det = b11 * b22 - b12 * b21
    cokernel = math.gcd(
        n * n,
        n * b11,
        n * b12,
        n * b21,
        n * b22,
        abs(det),
    )
    if cokernel <= 0 or (n * n) % cokernel:
        raise RuntimeError("Invalid lattice index")
    return (n * n) // cokernel


def powers(base: int, length: int, modulus: int) -> List[int]:
    vals = [1]
    b = base % modulus
    for _ in range(1, length):
        vals.append((vals[-1] * b) % modulus)
    return vals


def box_image_counts(M: int, A: int, p: int, N: int) -> Tuple[int, int]:
    p2 = powers(2, N, p)
    p3 = powers(3, N, p)
    pM = powers(M, N, p)
    pA = powers(A, N, p)

    row_points = {
        ((p2[a] * pM[b]) % p, (p3[a] * pA[b]) % p)
        for a in range(N)
        for b in range(N)
    }
    column_signatures = {
        ((p2[r] * p3[s]) % p, (pM[r] * pA[s]) % p)
        for r in range(N)
        for s in range(N)
    }
    return len(row_points), len(column_signatures)


def local_cyclic_slope(M: int, A: int, p: int) -> Tuple[bool, int | None, int]:
    """Return whether Q is in <P>, a least nonnegative slope, and ord(P)."""
    target = (M % p, A % p)
    x, y = 1, 1
    slope: int | None = None
    for k in range(p - 1):
        if slope is None and (x, y) == target:
            slope = k
        x = (x * 2) % p
        y = (y * 3) % p
        if (x, y) == (1, 1):
            order = k + 1
            return slope is not None, slope, order
    raise RuntimeError(f"Failed to find the order of P modulo {p}")


def slope_mismatch(M: int, A: int) -> float:
    return math.log(M) / math.log(2.0) - math.log(A) / math.log(3.0)


def analyze_pair(M: int, A: int, label: str, N: int, prime_limit: int) -> Tuple[List[dict], dict]:
    if M <= 0 or A <= 0 or N <= 0:
        raise ValueError("M, A, and N must be positive")

    rows: List[dict] = []
    R = N * N
    weighted_defect = 0.0
    primes_with_defect = 0
    primes_with_large_defect = 0
    cyclic_primes = 0
    rank_one_scale_primes = 0
    determinant_zero_primes = 0
    max_defect = 0

    for p in primes_upto(prime_limit):
        if (6 * M * A) % p == 0:
            continue
        g = primitive_root_prime(p)
        logs = discrete_log_table(p, g)
        matrix = (
            logs[2 % p],
            logs[M % p],
            logs[3 % p],
            logs[A % p],
        )
        det_integer = matrix[0] * matrix[3] - matrix[1] * matrix[2]
        det_mod = det_integer % (p - 1)
        full_image = image_size_mod_n(matrix, p - 1)
        row_count, col_count = box_image_counts(M, A, p, N)
        defect = max(R - row_count, R - col_count)
        cyclic, slope, orderP = local_cyclic_slope(M, A, p)

        weighted_defect += defect * math.log(p)
        primes_with_defect += int(defect > 0)
        primes_with_large_defect += int(defect >= R / 4)
        cyclic_primes += int(cyclic)
        rank_one_scale_primes += int(full_image <= p - 1)
        determinant_zero_primes += int(det_mod == 0)
        max_defect = max(max_defect, defect)

        rows.append(
            {
                "label": label,
                "M": M,
                "A": A,
                "p": p,
                "primitive_root": g,
                "b11": matrix[0],
                "b12": matrix[1],
                "b21": matrix[2],
                "b22": matrix[3],
                "det_mod_p_minus_1": det_mod,
                "full_orbit_image_size": full_image,
                "row_box_image_size": row_count,
                "column_box_image_size": col_count,
                "box_rank_defect_lower_bound": defect,
                "Q_in_cyclic_subgroup_P": int(cyclic),
                "least_nonnegative_local_slope": "" if slope is None else slope,
                "order_of_P": orderP,
            }
        )

    summary = {
        "label": label,
        "M": M,
        "A": A,
        "box_side_N": N,
        "box_cardinality": R,
        "prime_limit": prime_limit,
        "slope_mismatch_log2M_minus_log3A": slope_mismatch(M, A),
        "good_primes": len(rows),
        "primes_with_box_collision": primes_with_defect,
        "primes_with_defect_at_least_one_quarter": primes_with_large_defect,
        "primes_with_local_cyclicity": cyclic_primes,
        "primes_with_rank_one_scale_full_image": rank_one_scale_primes,
        "primes_with_discrete_log_determinant_zero": determinant_zero_primes,
        "maximum_box_defect_lower_bound": max_defect,
        "sum_defect_times_log_p": weighted_defect,
    }
    return rows, summary


def parse_pairs(specs: Sequence[str]) -> List[Pair]:
    pairs: List[Pair] = []
    for spec in specs:
        parts = spec.split(":")
        if len(parts) not in (2, 3):
            raise ValueError(f"Bad pair specification: {spec}")
        M, A = int(parts[0]), int(parts[1])
        label = parts[2] if len(parts) == 3 else f"M={M},A={A}"
        pairs.append((M, A, label))
    return pairs


def main() -> None:
    parser = argparse.ArgumentParser()
    parser.add_argument("--N", type=int, default=24)
    parser.add_argument("--prime-limit", type=int, default=5000)
    parser.add_argument(
        "--pair",
        action="append",
        default=[],
        help="M:A[:label]; may be repeated",
    )
    parser.add_argument("--csv", type=Path, required=True)
    parser.add_argument("--summary", type=Path, required=True)
    args = parser.parse_args()

    pairs = parse_pairs(args.pair) if args.pair else [
        (8, 27, "integral slope beta=3"),
        (5, 13, "small generic pair"),
        (3779, 467744, "extreme archimedean near-miss"),
        (9, 2, "synthetic coprime pair"),
    ]

    all_rows: List[dict] = []
    summaries: List[dict] = []
    for M, A, label in pairs:
        rows, summary = analyze_pair(M, A, label, args.N, args.prime_limit)
        all_rows.extend(rows)
        summaries.append(summary)

    args.csv.parent.mkdir(parents=True, exist_ok=True)
    with args.csv.open("w", newline="", encoding="utf-8") as fh:
        writer = csv.DictWriter(fh, fieldnames=list(all_rows[0].keys()))
        writer.writeheader()
        writer.writerows(all_rows)

    payload = {
        "summaries": summaries,
    }
    args.summary.write_text(json.dumps(payload, indent=2, sort_keys=True) + "\n", encoding="utf-8")
    print(json.dumps(payload, indent=2, sort_keys=True))


if __name__ == "__main__":
    main()
\end{lstlisting}

\begin{thebibliography}{99}

\bibitem{UnifiedReport}
\emph{Alaoglu--Erd\H{o}s Conjecture: Unified Progress Report},
internal research manuscript supplied with the present task, 2026.

\bibitem{AnthropicFable}
Anthropic,
\emph{Claude Fable 5}, public announcement, 11 August 2026.
\url{https://www.anthropic.com/news/claude-fable-5}.

\bibitem{AnthropicSystemCard}
Anthropic,
\emph{Claude Fable 5 System Card}, accompanying technical report, August 2026;
linked from the official Claude Fable 5 announcement.

\bibitem{AlpogeJacobian}
L.~Alp\"oge,
\emph{Counterexamples to the Jacobian conjecture},
arXiv:2608.02916, 2026.
\url{https://arxiv.org/abs/2608.02916}.

\bibitem{ByClaudeJacobian}
ByClaude,
\emph{Jacobian Conjecture Counterexamples}, project record, 2026.
\url{https://byclaude.ai/projects/jacobian-conjecture}.

\bibitem{ByClaudeHadamard}
ByClaude,
\emph{Hadamard Matrix Construction}, project record, 2026.
\url{https://byclaude.ai/projects/hadamard-matrix}.

\bibitem{VibeHadamard}
Vibe Mathed,
\emph{Hadamard Matrix Conjecture: order 668 and subsequent constructions},
public project notes, 2026.
\url{https://vibemathed.com/}.

\bibitem{VibeRank30}
Vibe Mathed,
\emph{Elliptic Curve Rank 30}, public project page describing an open target, accessed August 2026.
\url{https://vibemathed.com/}.

\bibitem{CorralesSchoof}
C.~Corrales-Rodrig\'a\~nez and R.~Schoof,
\emph{The support problem and its elliptic analogue},
J. Number Theory \textbf{64} (1997), 276--290.

\bibitem{PeruccaSupport}
A.~Perucca,
\emph{Two variants of the support problem for products of abelian varieties and tori},
J. Number Theory \textbf{129} (2009), 1883--1892;
arXiv:0712.2815.
\url{https://arxiv.org/abs/0712.2815}.

\bibitem{PeruccaDetect}
A.~Perucca,
\emph{On the problem of detecting linear dependence for products of abelian varieties and tori},
Acta Arith. \textbf{142} (2010), 119--128;
arXiv:0811.1495.
\url{https://arxiv.org/abs/0811.1495}.

\bibitem{PeruccaOrders}
A.~Perucca,
\emph{The order of the reductions of points on abelian varieties and tori},
Ph.D. thesis, Universiteit Leiden, 2008;
arXiv:0712.2812.
\url{https://arxiv.org/abs/0712.2812}.

\bibitem{PeruccaMultilinear}
A.~Perucca,
\emph{The multilinear support problem for products of abelian varieties and tori},
arXiv:1007.0455, 2010.
\url{https://arxiv.org/abs/1007.0455}.

\bibitem{BCZ2003}
Y.~Bugeaud, P.~Corvaja, and U.~Zannier,
\emph{An upper bound for the G.C.D. of $a^n-1$ and $b^n-1$},
Math. Z. \textbf{243} (2003), 79--84.

\bibitem{PinkSupport}
R.~Pink,
\emph{On the order of the reduction of a point on an abelian variety},
Math. Ann. \textbf{330} (2004), 275--291.

\bibitem{GelfondSchneider}
A.~O.~Gelfond,
\emph{Transcendental and Algebraic Numbers},
Dover, 1960;
and T.~Schneider, \emph{Einf\"uhrung in die transzendenten Zahlen}, Springer, 1957.

\end{thebibliography}

\end{document}
